\section{Closest Work and What Remains}\label{sec:related}
\textbf{scMINER itself} \cite{pan2025scminer}. MICA bins each cell's expression and scores each pair of cells with the distance $D = 1 - I/H$, where $I$ is their mutual information and $H$ their joint entropy. Up to 5,000 cells it embeds these distances with classical multidimensional scaling (MDS) and runs $k$-means. Above that it finds each cell's 80 nearest neighbors with HNSW, an approximate search index, embeds that graph with node2vec (random walks) and clusters the embedding with Louvain. On ten labeled sets it reports a mean ARI of 0.84, against 0.73 for Seurat and 0.68 for Scanpy. The labels chose every setting, and only PBMC14K is above 5,000 cells, so nine of the ten headline scores come from MICA given the true $k$. In the authors' own benchmark files MICA trails Scanpy on the three larger sets (76,000 to 647,000 cells), which the paper used only for run time and memory. Asked by a reviewer how to choose the resolution without labels, the authors suggested silhouette scores and marker genes but did not test them. The paper does not report how many true neighbors its HNSW graph finds, and no independent benchmark we know of includes MICA.

\textbf{Choosing the number of clusters without labels.} Yu et al. \cite{yu2022benchmarking} benchmarked 14 methods that estimate the number of cell types on mixtures of Tabula Muris cells. Most methods underestimated the count, more so as the true number grew, while SC3 and Seurat overestimated it, and a good ARI did not imply a correct count. MultiK \cite{liu2021multik} reruns Seurat on many 80\% subsamples and proposes the cluster counts that recur most stably, and chooseR \cite{pattersoncross2021chooser} picks the finest resolution whose clusters still hold together across subsamples. A 2025 preprint \cite{szmigiel2025benchmarking} reports that on six of 16 datasets the Scanpy default did no better than a typical random setting. None of these includes a mutual information method, and the two selection tools were shown only with Seurat or Scanpy pipelines. Ullmann et al. \cite{ullmann2023overoptimistic} showed that tuning a new clustering algorithm on the ground truth overstates what a user without labels will get, and our main question applies that warning to scMINER.

\textbf{Testing whether clusters are real.} sc-SHC \cite{grabski2023significance} tests each split of a hierarchy against a one-population null model and controls the family-wise error. On data from a single cell line, Seurat's Louvain and Leiden at their defaults both returned 9 clusters, while sc-SHC returned 1. CHOIR \cite{sant2025choir} first overclusters, then merges neighboring clusters that a random forest cannot tell apart better than chance. On 230 simulated sets it matched or beat 15 other methods. Neither was tested on a mutual information method, but each gives a label-free stopping rule that can be applied to MICA and Leiden alike.

\textbf{Neighbor graphs and hubness.} Proposing candidates cheaply and then rescoring them exactly is a standard pattern in nearest neighbor search, as in the refinement step of Faiss \cite{douze2024faiss}, and MICA itself uses HNSW \cite{malkov2020efficient}. Song et al. \cite{song2012comparison} found that binned mutual information closely tracks correlation for gene pairs, which points to correlation as the cheap step. Amblard et al. \cite{amblard2022hubness} found that single-cell neighbor graphs have hub cells, which appear in far more neighbor lists than a typical cell. Correcting for hubs improved Leiden clustering, though with the resolution tuned to the true number of clusters. We know of no study of the neighbors a mutual information graph misses, or of its hubs.

\textbf{Annotation with language models.} GPTCelltype \cite{hou2024gpt4} asks GPT-4 to name clusters from their top 10 marker genes, and the names fully or partly matched manual ones for over 75\% of cell types in most tissues. CASSIA \cite{xie2025cassia} adds agents that check the cited markers, and GeneAgent \cite{wang2025geneagent} checks each claim about a gene set against 18 databases. SCENIC \cite{aibar2017scenic} and metaVIPER \cite{ding2018metaviper} infer regulator activity in single cells, but no annotator we found uses it as evidence, and only GeneAgent checks claims one by one.

\textbf{What remains.} As far as we can tell, no published work compares mutual information clustering with a matched PCA and Leiden baseline under label-free selection, or checks MICA's graph against an exact answer. Our project does both.
